\section*{Глава 4. Оформление итогов работы}
\addcontentsline{toc}{section}{Глава 4. Оформление итогов работы}
\stepcounter{section}

\subsection{Создание Git-репозитория с кодом проекта}

Код проекта размещён в системе контроля версий Git и опубликован в
публичном репозитории на платформе GitHub. Перед публикацией
репозиторий был очищен от файлов, не подлежащих хранению в системе
контроля версий: служебных файлов IDE (\texttt{.idea}), файлов
операционной системы (\texttt{.DS\_Store}), а также файлов окружения
(\texttt{.env}), содержащих секретные параметры (ключи, пароли БД) и
каталога зависимостей \texttt{node\_modules}. Правила исключения
оформлены в файле \texttt{.gitignore} (листинг~\ref{lst:gitignore}),
заданном как на уровне всего репозитория, так и отдельно для
серверной и клиентской части проекта.

\begin{lstlisting}[language=, caption={Правила исключения файлов из репозитория (.gitignore)}, label={lst:gitignore}]
# OS
.DS_Store

# IDE
.idea/

# Python
venv/
__pycache__/
db.sqlite3

# Env / secrets
.env
.env.*
!.env.example

# Frontend
node_modules/
dist/
\end{lstlisting}

Файлы, ранее случайно попавшие в репозиторий (\texttt{.idea},
\texttt{.DS\_Store}), были удалены из индекса Git командой
\texttt{git rm -{}-cached} без удаления их с диска разработчика.
Секретные параметры окружения (пароль базы данных, ключ DeepSeek API,
секретный ключ Django) хранятся исключительно в файле \texttt{.env} на
сервере и никогда не передаются через систему контроля версий; для
разработчиков, разворачивающих проект впервые, в репозитории
предусмотрен шаблон \texttt{.env.example} с описанием требуемых
переменных без реальных значений.

\subsection{Деплой приложения на хостинг}

Для развёртывания приложения выбран облачный сервер (VPS) с
конфигурацией 2 vCPU, 4~ГБ оперативной памяти, 40~ГБ NVMe-накопителя и
пропускной способностью сети 1~Гбит/с под управлением ОС Ubuntu~24.04
LTS. Выбор облачного сервера, а не платформ типа PaaS (Render,
Railway), обусловлен требованиями подсистемы проверки решений Judge0
CE к привилегированному режиму контейнеров (\texttt{privileged: true})
для работы изолятора \texttt{isolate}, который большинство PaaS-провайдеров
не поддерживает.

Развёртывание выполняется единой командой \texttt{docker compose up -d
-{}-build}, поднимающей все восемь контейнеров системы (см.
листинг~\ref{lst:compose} главы~3) в едином окружении. Для всех
сервисов задана политика перезапуска \texttt{restart: unless-stopped},
обеспечивающая автоматическое восстановление контейнеров после сбоя
или перезагрузки сервера.

При подготовке к промышленной эксплуатации устранён ряд проблем
конфигурации, выявленных в процессе развёртывания:

\begin{itemize}
	\item значение \texttt{ALLOWED\_HOSTS} Django дополнено внутренним
	      именем сервиса \texttt{backend}, поскольку HTTP-заголовок
	      \texttt{Host} запросов, проксируемых сервером разработки Vite
	      внутри Docker-сети, отличается от внешнего IP-адреса сервера;
	\item публичные порты СУБД PostgreSQL (5432) и сервиса Judge0 CE
	      (2358) закрыты в файле \texttt{docker-compose.yml}: оба сервиса
	      используются только внутри внутренней docker-сети и не должны
	      быть доступны из внешней сети, поскольку Judge0 CE без
	      дополнительной настройки ключа авторизации позволяет выполнять
	      произвольный код, а открытый порт СУБД создаёт риск
	      несанкционированного доступа к данным пользователей;
	\item пароль СУБД приложения переведён на единый источник
	      (\texttt{env\_file: ./backend/.env}), используемый одновременно
	      контейнером \texttt{db} и серверным приложением, что исключило
	      рассинхронизацию значения пароля между файлами конфигурации.
\end{itemize}

Для автоматизации обновления приложения после внесения изменений в
код настроен процесс непрерывного развёртывания (CI/CD) на основе
GitHub Actions. При каждом помещении (\textit{push}) изменений в ветку
\texttt{main} запускается сценарий (листинг~\ref{lst:cicd}), который по
протоколу SSH подключается к серверу, загружает обновлённый код
(\texttt{git pull}) и пересобирает изменившиеся контейнеры. Учётные
данные для подключения (приватный SSH-ключ, адрес и имя пользователя
сервера) хранятся в защищённых секретах репозитория GitHub Actions и не
фигурируют в коде проекта.

\begin{lstlisting}[language=, caption={Сценарий автоматического развёртывания (.github/workflows/deploy.yml)}, label={lst:cicd}]
name: Deploy

on:
  push:
    branches: [main]

jobs:
  deploy:
    runs-on: ubuntu-latest
    steps:
      - name: Deploy to server
        uses: appleboy/ssh-action@v1.0.3
        with:
          host: ${{ secrets.SSH_HOST }}
          username: ${{ secrets.SSH_USER }}
          key: ${{ secrets.SSH_PRIVATE_KEY }}
          script: |
            cd ~/algo/algo_platform
            git pull
            docker compose up -d --build
\end{lstlisting}

Таким образом, после внесения изменений в код и их публикации в
основной ветке репозитория обновлённая версия приложения автоматически
становится доступна на сервере без необходимости вручную подключаться
к нему и выполнять команды развёртывания.

\subsection{Оформление отчёта о проделанной работе}

Пояснительная записка к курсовой работе подготовлена с использованием
системы компьютерной вёрстки LaTeX на основе шаблона, утверждённого
кафедрой «Инфокогнитивные технологии» Московского Политеха, и включает
титульный лист, задание на курсовую работу, содержание, введение,
основную часть (главы 1--4), заключение, список использованных
источников и приложения с вынесенными диаграммами. Структура отчёта
соответствует перечню вопросов, подлежащих разработке, указанному в
задании на курсовую работу (табл.~1 задания): главы пояснительной
записки последовательно раскрывают анализ предметной области,
проектирование, разработку и оформление итогов работы.

Листинги исходного кода, приведённые в главе 3, оформлены средствами
пакета \texttt{listings} с указанием пути к соответствующему файлу
проекта в подписи, что позволяет однозначно сопоставить приведённый в
отчёте код с актуальной версией, опубликованной в Git-репозитории
проекта.
